% commentaire.tex — f163r, vv. 43–79 (conversion initiale des travaux antérieurs)
\comm{43–7}{ἐχέτω δὲ πνοὰς νήνεμος αἰθήρ ...}{De même que dans les \italique{Bacchantes} d'Euripide (v. 1084), à la voix de Dionysos « l'éther se tut, le ravin herbeux étouffa le murmure des feuilles, et l'on n'eût point entendu le cri des bêtes sauvages ».\citComm{σίγησε δ᾽ αἰθήρ, σῖγα δ᾽ ὕλιμος νάπη\\
φύλλ᾽ εἶχε, θηρῶν δ᾽ οὐκ ἂν ἤκουσας βοήν.}Et aussi dans Théocrite (II, 38), tout se tait à l'arrivée d'Hécate. (Note de Willems.)}
\comm{43}{ἐχέτω}{est mis pour ἐπεχέτω (ici : \italique{retenir}) comme l'indique la scholie intérieure au texte du \italique{Ravennas}.}
\comm{45}{βομβάξ.}{Forgée sur le modèle d'ὄμπαξ -- interjection pour réclamer le silence « paix ! (πάξ) », « assez ! » --, βομβάξ fait entendre le \italique{bourdonnement}, βόμβος (ὁ βομβυλιός désigne tout type d'insecte bourdonnant). Plaisante manière pour Aristophane de ridiculiser l'image détournée du serviteur : le bruyant thiase et sa pompe de muses venues honorer la demeure d'Agathon sont ramenées à la plus commune métaphore de la \italique{ruche}.}
\comm{49–50}{μέλλει γὰρ ὁ καλλιεπὴς Ἀγάθων / πράμος ἡμέτερος.}{πράμος n'est ni autrement attesté ni même relevé par les grammairiens et doit nécessairement être scandé $\cup \cup $ pour rentrer dans le mètre anapestique, empêchant tout rapprochement étymologique pertinent. 
Autant d'indices qui font soupçonner une erreur textuelle plutôt qu'un hapax agathonien.
πρόμος, \italique{le chef}, fut proposé par Grynaeus et quasi-unanimement repris depuis. 
La correction est très légère mais n'est pas satisfaisante : 
il manque au terme πρόμος une acception \italique{sensu obsceno} pour expliquer l'exclamation de Mnésiloque : μῶν βινεῖσθαι;\nc{Sommerstein voit en πρόμος une légère déformation suggérant à l'oreille πόρνος... « Notre chef de file --- Celui qu'on enfile ? » Mais que c'est loin en grec !.. Certains ont voulu justifier le rapprochement πρόμος -- πόρνος en arguant de la plaisanterie de \italique{Nuées} (v. 394) sur l'indissociabilité entre βροντή et πορδή -- certes pas au niveau de la prononciation. Nous avons profité des annexes auxquelles nous renvoyons le lecteur pour exégéter ce jeu de mots qui roule sur la langue proverbiale.} 
De plus, l'on s'attendrait plutôt à trouver un infinitif réclamé par μέλλει. Une glose d'Hésychius va nous permettre de débrouiller l'écheveau. Cette glose est elle-même légèrement corrompue (π 3188) :\citComm{†πραῖνοι†· πρηνίζειν, καταστρέφειν.}Sans nous attarder pour l'instant sur les hypothèses de corruptions, constatons que Hésychius atteste l'existence du verbe πρᾳνόω (ou πρανόω\nc{Même genre de variantes orthographiques entre πρᾷος et πρᾶος.}).
Pour le sens, πρηνίζειν signifie \italique{précipiter}, \italique{renverser}, \italique{jeter} (en avant).
Le serviteur d'Agathon construit ainsi sa phrase :\citComm{μέλλει γὰρ Ἁγάθων πρανοῦν δρυόχους\nc{Les δρύοχοι sont les \italique{chantiers}, ou poutres sur lesquelles s'établit la quille d'un vaisseau en construction. On rapprochera la métaphore de \italique{Guêpes} 30 : τὴν τρόπιν τοῦ πράγματος, « la base de l'affaire ». (Note de Taillardat §758 n. 1 ; voir aussi §673.)}.\\
Agathon s'apprête à jeter les fondations.}Seulement, πρηνίζειν peut désigner aussi l'action de s'accroupir tête baissée\nc{Cf. Hésychius π 3199, π 3269 πρηνιχθέντα· ἐπὶ πρόσωπον καὶ ἐπὶ στόμα πεσόντα ~ Photius π 1138 et L'Etymologicum genuinum : πρηνίξαι πρηνιχθῆναι καὶ πεπρηνίχθαι· γνὺξ γάρ ἐστιν κυρίως τὸ εἰς γόνυ πεσεῖν ... τὸ δὲ πρηνιχθῆναι ἐπὶ τὸ πρόσωπον [πεσεῖν] et d'autres références dans l'abondant et récent apparat des témoins de l'édition Theodoridis de Photius.}. (πρανοῦν, comme accroupir, fait entendre la \italique{croupe}, πρανής.) C'est bien-sûr dans ce sens détourné que le prend Mnésiloque -- image d'un Agathon se jetant à son travail dans la position du penseur de Rodin ? L'occasion est trop belle : Mnésiloque s'imagine Agathon en position κύβδα (cf. le vers 489) et tournant le dos à d'éventuels assaillants --- μῶν βινεῖσθαι; demande-t-il. Le véritable Agathon devait effectivement travailler dans une posture qui prêtait à rire aux mauvais esprits : c'est sur son lit et dans cette même posture qu'il apparaît sur scène quelques vers plus loin ; cette position étrange lui vaudra d'être comparé par Mnésiloque à Cyrène \italique{aux douze postures} (δωδεκαμήχανος) (v. 98).\\
Reste à expliquer la présence dans le \italique{Ravennas} de †ἡμέτερος† après †πράμος†. Hésychius glose †πραῖνοι† par deux infinitifs, πρηνίζειν et καταστρέφειν, preuve que le mot relevé en vedette devait originellement être lui aussi un infinitif. Le rétablissement de πρανοῦν se fait alors sans difficulté. C'est la même mécoupure du verbe πρανοῦν, lu †πρανοι \ ν† qui pourrait expliquer le †πράμος† et le début de †ἡμέτερος† du \italique{Ravennas} selon le \italique{stemma} suivant :  \begin{align*}
\overset{\overset{\text{\footnotesize{\allemand{Urtext}}}}{\text{πρηνίζειν, καταστρέφειν}}}{\text{πρανοῦν [...]}}\xrightarrow{\text{mécoup. \ } + \text{\ conf. \  υ/ι}}\overset{\overset{\text{\footnotesize{Archétype}}}{\text{πρηνίζειν, καταστρέφειν}}}{\text{πρανοι \ \ ν[...]}}&\xrightarrow{\text{\ \ \ \ \ }}\overset{\text{\footnotesize{Hésychius π 3188}}}{\text{πραῖνοι· πρηνίζειν, καταστρέφειν}}\\
& \xrightarrow{\text{confus. \ ν/μ et ι/ς}}\overset{\text{\footnotesize{\italique{Ravennas}}}}{\text{πραμος <η>μ[ετερος]}}
\end{align*}
Rien de plus semblable qu'un ι et un ς, ainsi qu'un μ et un ν  -- dans les manuscrits écrits en minuscules à l'instar du \italique{Ravennas}.
Il s'agit maintenant de comprendre la présence de †ετερος†. Sans tenter d'explications, je constate qu'au vers 562 des \italique{Nuées}, les deux principaux manuscrits des pièces d'Aristophane, notre \italique{Ravennas} et le \italique{Venetus}, portent deux leçons différentes : ἑτέρας \italique{versus} προτέρας\nc{Les écrits anciens comportaient de nombreuses abréviations. (Cf. par exemple la très ingénieuse reconstruction du vers 80 par Nauck.) Par exemple, προσ était souvent abrégé d'une sorte de Ʒ renversé qui ressemblait à un epsilon : \renverse{Ʒ}. Voir G. Bady, \italique{Paléographie grecque : éléments de cours}, ou bien V. Gardthausen, \italique{Griechische palaeographie}\expo{2}, II (1913), p. 340.}. 
Je suppose ici la même corruption :\citComm{μέλλει γὰρ ὁ καλλιεπὴς Ἀγάθων πρανοῦν προτέρους δρυόχους, τιθέναι δράματος ἀρχάς. κάμπτει δὲ νέας ἁψῖδας ἐπῶν κτλ.\\
Agathon au beau langage s'apprête à jeter les premiers fondements, à disposer les prémices d'un drame. Il tourne de nouvelles roues de mots etc.}}
\comm{52–7}{τιθέναι δράματος ἀρχάς... καὶ χοανεύει.}{Agathon le tragique n'est pas un génie bâtisseur, un τέκτων ; il ne crée pas : il fabrique, il bricole. Il est donc tout-à-tour charpentier de marine, menuisier, modeleur : « il va poser des chantiers, fondements d'un drame. etc. » Ces métaphores sont évidemment burlesques ; elles avilissent délibérément l'art d'Agathon, ravalé au rang de métier et ne peuvent que tourner le poète en dérision...\\
Agathon est donc à mille lieues du mystère de l'inspiration, décrit par Platon (dans l'\italique{Ion}) : κοῦφον γὰρ χρῆμα ποιητής ἐστιν καὶ πτηνὸν καὶ ἱερόν, « le poète est chose légère, ailée, sacrée ». De fait, Platon connaît bien les images du menuisier et du modeleur\nc{Voir Pierre Louis, \italique{Les Métaphores de Platon}, Belles Lettres, 1945, p. 98, 170, 204-5.} mais jamais il ne les applique au poète (ἀποτορνεύειν est dit de Lysias, \italique{Phèdre} 243e : ἀκριβῶς ἕκαστα τῶν ὀνομάτων ἀποτετόρνευται et reste péjoratif. La même observation vaut pour κολλᾶν, en parlant d'un auteur, \italique{Ménexène} 263b ; \italique{Phèdre} 278e). À noter que ces images semblent banales ; ainsi, dans le \italique{Timée} (81b), on trouve : οἷον ἐκ δρυόχων, « comme depuis les longerons de la carcasse », c'est-à-dire \italique{entièrement}. κάμπτειν : \italique{courber}, \italique{plier} désigne en musique les courbes capricieuses d'un air ou les inflexions de la voix. (Notes de Taillardat\nc{\italique{Les images d'Aristophane}, §758 et §784.}.)}
\comm{56–7}{καὶ κηροχυτεῖ καὶ γογγύλίζει / καὶ χοανεύει — καὶ λαικάζει.}{La métaphore du serviteur est tirée de la technique de la \italique{cire perdue}\nc{Voir sur cette technique l'article de C. C. Mattusch « The Berlin Foundry Cup: The Casting of Greek Bronze Statuary in the Early Fifth Century B. C. », \italique{American Journal of Archaeology}, 84 n° 4 (1980), pp. 435-444, ou son livre \italique{Greek Bronze Statuary: From the Beginnings Through the Fifth Century B.C.} (1988).} ; un procédé de moulage pour obtenir une statue en métal à partir d'un modèle en cire : on commence par modeler à la cire la statue (κηροχυτεῖν), modèle qu'on recouvre d'une gangue d'argile qui fera office de moule. C'est à cette opération que renvoie probablement le verbe γογγυλίζειν (corr. γογγύλλειν), associant la gangue à une \italique{rave ronde}, γογγυλίς. Il faut laisser à ce moule d'argile des cheminées, certaines pour insérer des entonnoirs et y couler le métal en fusion (χοανεύειν) qui va prendre la place de la cire, et d'autres cheminées pour permettre à la cire fondue d'être évacuée.
Cet esprit tordu de Mnésiloque y voit la description d'une scène érotique.
Il interrompt d'une obscénité: « καὶ λαικάζει ».
Ce verbe a le sens de \latin{fellare}\nc{\italique{Confer} par exemple Chantraine, \italique{Dictionnaire étymologique de la langue grecque}. γογγύλλειν peut avoir le sens de \italique{rendre compact}, \italique{serré}. Cf. la \nieme{3} note §166 de Taillardat : \italique{ramassé}, \italique{compact}, \italique{serré} sont les sens que Taillardat donne à l'adjectif γογγύλος aux vers 25-28 de la \italique{Paix}.}.}
\comm{61}{γογγυλίσας καὶ συστρέψας.}{Vers retouché \italique{gratia metri} d'après un passage de \italique{Lysistrata} (v. 975), jouée probablement l'année précédente (cf. Introduction).
Allons-y voir d'un peu plus près pour comprendre précisément l'image de Mnésiloque : Kinésias est dans un état lamentable après que Myrrhine l'ait fortement échaudé... et finalement laissé raide. Le demi-chœur des hommes en vient à souhaiter un traitement de faveur pour cette s[célérate] (v. 972-979) : \citComm{μιαρὰ δῆτ' ὦ Ζεῦ Ζεῦ· |
εἴθ' αὐτὴν ὥσπερ τοὺς θωμοὺς |
μεγάλῳ τυφῷ καὶ πρηστῆρι |
\textbf{ξυστρέψας καὶ ξυγγογγύλας} |
οἴχοιο φέρων, εἶτα μεθείης, |
ἡ δὲ φέροιτ' αὖ πάλιν ἐς τὴν γῆν, |
κᾆτ' ἐξαίφνης |
περὶ τὴν ψωλὴν περιβαίη.\\
Une gueuse ! ô Zeus, Zeus, puisse-t-elle, comme les tas de blé dans un coup de vent et une tourmente, roulée par toi [Kinésias] et ramassée en boule, être portée dans l'air, puis, lâchée, retomber à terre et du coup enfourcher ce priape.\\
\textbf{975} ξυγγογγυλίσας R, Γ, Suid. : \lecon{ξυγγογγύλας} Cobet}L'image est tirée de la technique du foulage et du vannage du grain. Voici ce qu'en dit le Daremberg \& Saglio\nc{\italique{Dictionnaire des antiquités grecques et romaines}, IV-2, pp. 906-907. Article « \latin{Rustica Res} ».} :\citComm{Foulage (ἀλόησις) --- Les Grecs ne battaient les céréales ni en grange, ni en plein air, car ils n'avaient point de fléau pour faire sortir le grain de l'épi. Ils employaient la vieille méthode égyptienne du dépiquage ou foulage, qui n'est plus guère en usage que chez les Provençaux et les Languedociens. En juin, c’est-à-dire, à l'apparition d’Orion\nc{Hésiode, \italique{Travaux et les jours}, v. 596.}, on portait les gerbes sur une aire, qualifiée de μεγάλη\nc{\italique{Iliade}, XIII, v. 588.} , ἐϋκτίμενη\nc{\italique{Iliade}, XX, v. 496.}, ἐϋτροχάλος\nc{Hésiode \italique{Travaux et les jours}, v. 597.}, ἱερά\nc{\italique{Iliade}, V, v. 499.}, et construite, plus ou moins, d'après le modèle de l'aire sacrée d'Éleusis\nc{Pausanias, I, 38.}. Des ouvriers, ἐπαλώστες\nc{Xénophon, \italique{Économiq.}, XVIII.}, dressent les gerbes a côté les unes des autres en les inclinant un peu vers le centre et on fait entrer les animaux : bœufs, chevaux, mulets, parfois même des ânes\nc{\italique{Anthol. palat.} IX, 301.} ; ils sont accouplés deux à deux et marchent en cercle, d’abord au pas, puis au trot, sous la direction d'un homme, placé au centre de l'aire, qui tient en main les guides, δεσμά\nc{\italique{Anthol. palat.} VI, 104.}. D'autres ouvriers, armés de fourches en bois, τρίναξ\nc{\italique{Anthol. palat.} VI, 95.}, régularisent la foulaison, ὁμαλιεῖται ὁ ἁλοητός, en retournant la paille et en la ramenant sous les pieds des animaux, pour faciliter le manège, δῖνος.\\
Nettoyage du grain (λίκμησις)  — Les Grecs ne connaissaient point le vannage proprement dit, opération de nos batteurs en grange, qui nettoye simplement le grain des glumes ou glumelles et que l'on exécute près de la porte d’une grange, dans un léger courant d'air, mais à l'abri du vent. Au contraire, les Grecs attendaient que le vent fût assez fort\nc{\italique{Iliade}, V, vv. 499 sq. ; XIII, vv. 588 sq. ; cf. \italique{Anth. palat.}, VI, 53.} pour emporter les fétus de paille et leur permît de pratiquer le \italique{ventage}, qui peut se faire de deux façons :  1°  Avec une corbeille, λικμός, πλόκανον [\latin{Ventilabrum}] [...]    2°   Avec une pelle en bois, πτύον [\latin{Pala}]. C'est notre \italique{nettoyage à la roue} ; on jette circulairement contre le vent et, aussi loin que possible, à la hauteur d’un mètre au moins, une pelletée de grains qui retombent sur le sol, tandis que la paille et les balles sont emportées à l'extrémité de l'aire.}Ces longues mais utiles précisions ayant été apportées, le passage de \italique{Lysistrata} est maintenant clair : Myrrhine doit être traitée comme un fétu de paille qu'on retourne et compactifie en boule (ξυστρέψας καὶ ξυγγογγύλας\nc{Soit dit en passant, c'est certainement l'un de ces deux passages que glose Hésychius (γ 764) : γογγύλλειν· συστρέφειν.}) avant de le jeter (εἶτα μεθείης) pour qu'il jonche l'aire vide ; ici ἡ ψωλή $=$ τὸ κενόν, mais avec un double sens obscène : le [phallos] \italique{nu}, c'est-à-dire décalotté.\\
Revenons aux \italique{Thesmophories}. Le serviteur demande qui est le paysan (τίς ἀγροιώτας) qui se permet d'interrompre de la sorte son poème. Mnésiloque lui répond donc par une image agricole certainement populaire : il se prépare à le retourner et le mettre en boule (συγγογγύλας καὶ συστρέψας v. 61) aussi rustiquement que le chœur des vieillards voulait faire avec Myrrhine.}
\comm{67–8}{χειμῶνος οὖν / ὄντος κατακάμπτειν τὰς στροφὰς οὐ ῥᾴδιον.}{Les Thesmophories avaient lieu du 11 au 13 du mois de Pyanepsion, jours qui tombent en général pendant notre seconde moitié du mois d'octobre.
Le serviteur continue de filer sa métaphore : la cire est plus facile à modeler lorsqu'elle est exposée au Soleil.}
\comm{70–4}{τί οὖν ἐγὼ δρῶ; ... ὄντα κηδεστὴν ἐμόν.}{La répartition des ce passage diffère d'un éditeur à un autre ; chacun voulant apporter son grain de sel. Inscrivons-nous, nous aussi, dans cette grande tradition philologique et apportons notre pierre à l'édifice : une nouvelle répartition des répliques qui, espérons-le, clôturera un débat qui n'a que trop duré. La répartition des deux répliques du vers 70 ne pose aucun problème : il suffit de songer aux caractères des trois personnages en présence : le ton empoté du « τί οὖν ἐγὼ δρῶ; » ne peut convenir qu'au personnage pour l'instant passif et qui subit l'action : Mnésiloque\nc{Réplique analogue de Mnésiloque au vers 925 : ἐγὼ δ᾽ ὁ κακοδαίμων τί δρῶ;} (Euripide ruisselle de défauts mais garde toujours une certaine dignité ; jamais il ne s'abaisserait à lancer un mot aussi trivial et décontenancé – ses plaintes ont davantage de lyrisme !) La réponse : « περίμεν᾽, ὡς ἐξερχεται » ne peut être attribuée qu'à Euripide : c'est lui qui a besoin de Mnésiloque ; de plus, le serviteur ne délivre ni conseil ni ordre : il se contente de répondre le plus sèchement possible aux questions (cette répartition est d'ailleurs celle du manuscrit de Ravenne). Les vers 72 à 74 sont évidemment prononcés par Mnésiloque qui cherche à percer la détresse d'Euripide et nous révèle par la même occasion son lien de parenté avec le poète. Reste le vers 71 : il est plus que tentant de l'attribuer à Mnésiloque et faire ainsi des quatre vers 71-74 l'exact pendant des quatre vers 1-4 : une invocation à Zeus suivie d'une question : là pour savoir où Euripide le mène et ici pour savoir ce qui le tourmente. Problème : la présence d'un \italique{signe} (σημεῖον), probablement un \italique{paragraphos}, sur d'anciens manuscrits a été interprétée comme indiquant un changement d'interlocuteurs entre le vers 71 et les vers 72-74. Mais tel n'est pas le cas ici, et c'est ce que nous indiquait une ancienne scholie, très précieuse et incomprise\nc{Incomprise déjà de Suidas – lexicographe hypothétique du \nieme{IX} ou \nieme{X} siècle, à qui l'on attribue le lexique de la \italique{Souda}. Ce dernier fourmille en références et citations d'Aristophane : la \italique{Souda} est un outil de recherche sur Aristophane indispensable. 
Il en existe une version numérique : \href{https://www.cs.uky.edu/~raphael/sol/sol-html/}{https://www.cs.uky.edu/~raphael/sol/sol-html/} – incomprise donc de Suidas ce qui tend à prouver son ancienneté : Suidas (κ 1487) biffe le lemme `τὸ σημεῖον', pourtant indispensable à la bonne compréhension et égare ainsi les différents éditeurs des scholies qui l'appliquent au vers 74, voire qui la retouchent pour finir de la rendre tout-à-fait inintelligible.} : 
\citComm{τὸ σημεῖον· οὐ μόνον κηδεστὴς ὁ ἐκδιδούς, ἀλλὰ καὶ ὁ λαμβάνων.\\
\francais{Le signe [typographique] : non seulement le Parent est celui qui abandonne [la parole], mais aussi celui qui la [re]prend.}}La scholie ne précise pas à quoi servait ce σημεῖον ; serait-ce une pause ? quoi qu'il en soit, le protagoniste devait certainement faire quelque jeu de théâtre après l'invocation à Zeus du vers 71, identique à celui déjà fait après le premier vers et montrer ainsi aux spectateurs le parallélisme entre les deux quatrains.}
\comm{71}{ὦ Ζεῦ τί δρᾶσαι διανοεῖ με τήμερον;}{Formule courante ; se reporter, par exemple, aux \italique{Oiseaux}, vv. 58 \& 62.}
\comm{75}{προπεφυραμένον}{Pour comprendre cette image, l'on se rapportera à un passage des \italique{Oiseaux} (v. 462-463) et le commentaire qu'en fait Taillardat (§755) :
\citComm{καὶ προπεφύραται λόγος εἷς μοι, / ὃν διαμάττειν οὐ κωλύει\\
\francais{j'ai sous la main un discours déjà bien délayé et rien n'empêche que j'achève de le pétrir.}}« Φυρᾶν signifie proprement \italique{délayer} de la farine dans un liquide\nc{Cf. Hippocrate \italique{Du régime} (\italique{Περί διαίτης}) 2,40 et Thucydide 3,49.}. » Est donc προ-πεφυραμένον ce qui a déjà été délayé à l'avance mais point encore pétri ; autrement dit : le complot est ourdi mais rien n'est encore joué. L'avantage de ce genre de métaphore boulangère est qu'elle se transpose assez aisément en français\nc{Comme l'épisode raconté par Démosthène dans les \italique{Cavaliers} (v. 55-57) : il s'est fait voler par Cléon les pains qu'il avait distribués... alors que c'est lui qui avait mis la main à la pâte ! Le verbe \italique{pétrir} et le \italique{pétrin} ont donné nombre d'expressions en français.}.}
\comm{76–7}{τῇδε θἠμέρᾳ κριθήσεται / εἴτ᾽ ἔστ᾽ ἔτι ζῶν εἴτ᾽ ἀπόλωλ᾽ Εὐριπίδης.}{Le style est paratragique (avec l'expression poétique `τῇδε θἠμέρᾳ' plutôt que la forme plus courante τήμερον). Apprécier surtout à sa juste valeur l'allitération ! Lyrisme ? emphase ? Les `ttttttttt' saccadent-ils le vers ou lui donnent-ils du rythme ? sans doute un peu des deux. (Euripide est souvent à la frontière du sublime et du ridicule, frontière qu'il passe et repasse allègrement).}
